\documentclass[12pt,a4paper]{article}
\usepackage[utf8]{inputenc}
\usepackage[T1]{fontenc}
\usepackage{amsmath, amssymb}
\usepackage{enumitem}
\usepackage[margin=2cm]{geometry}

\begin{document}

\begin{center}
    \textbf{\Large Latihan Soal IPAS SMK Kelas X} \\
    \textbf{Persiapan Ulangan Akhir Semester}
\end{center}
\vspace{0.5cm}

\begin{enumerate}[leftmargin=*]

\item Tempat di permukaan bumi, baik secara keseluruhan maupun hanya sebagian, yang digunakan oleh makhluk hidup untuk tinggal disebut \dots
    \begin{enumerate}[label=\alph*.]
        \item Ruang
        \item Peta
        \item Atlas
        \item Globe
        \item Wilayah
    \end{enumerate}

\item Gambaran sebagian atau seluruh permukaan bumi pada suatu bidang datar yang diperkecil dengan menggunakan skala tertentu dinamakan \dots
    \begin{enumerate}[label=\alph*.]
        \item Globe
        \item Atlas
        \item Peta
        \item Proyeksi
        \item Topografi
    \end{enumerate}

\item Peta yang menggambarkan kenampakan bumi secara keseluruhan, baik kenampakan alam maupun buatan manusia disebut \dots
    \begin{enumerate}[label=\alph*.]
        \item Peta khusus
        \item Peta umum
        \item Peta tematik
        \item Peta tata guna lahan
        \item Peta geologi
    \end{enumerate}

\item Berikut ini yang merupakan contoh dari peta tematik atau peta khusus adalah \dots
    \begin{enumerate}[label=\alph*.]
        \item Peta topografi
        \item Peta korografi
        \item Peta kepadatan penduduk
        \item Peta dunia
        \item Peta benua Asia
    \end{enumerate}

\item Sekumpulan peta yang dirangkai, disusun, dan dibukukan menjadi satu kesatuan disebut \dots
    \begin{enumerate}[label=\alph*.]
        \item Ensiklopedi
        \item Katalog
        \item Buku teks
        \item Atlas
        \item Globe
    \end{enumerate}

\item Model tiruan bola bumi dalam bentuk tiga dimensi yang dibuat sedemikian rupa sehingga paling mendekati bentuk bumi sesungguhnya adalah \dots
    \begin{enumerate}[label=\alph*.]
        \item Peta relief
        \item Maket
        \item Miniatur
        \item Globe
        \item Atlas
    \end{enumerate}

\item Keunggulan mutlak globe dibandingkan dengan peta datar adalah \dots
    \begin{enumerate}[label=\alph*.]
        \item Mudah dibawa kemana-mana
        \item Harganya lebih murah
        \item Meminimalisir distorsi bentuk, luas, jarak, dan arah
        \item Dapat dilipat dengan mudah
        \item Menampilkan data secara lebih detail di satu lembar
    \end{enumerate}

\item Simbol pada peta yang digunakan untuk menyajikan tempat atau data posisional, seperti ibu kota negara dan gunung berapi, adalah simbol \dots
    \begin{enumerate}[label=\alph*.]
        \item Garis
        \item Warna
        \item Area
        \item Titik
        \item Luasan
    \end{enumerate}

\item Pada peta, jalan raya dan rel kereta api biasanya digambarkan dengan menggunakan simbol \dots
    \begin{enumerate}[label=\alph*.]
        \item Titik
        \item Area
        \item Garis
        \item Warna
        \item Huruf
    \end{enumerate}

\item Warna biru pada peta umumnya digunakan untuk mewakili \dots
    \begin{enumerate}[label=\alph*.]
        \item Dataran rendah
        \item Dataran tinggi
        \item Pegunungan
        \item Wilayah perairan
        \item Hutan
    \end{enumerate}

\item Angka perbandingan antara jarak pada peta dengan jarak sebenarnya di permukaan bumi disebut \dots
    \begin{enumerate}[label=\alph*.]
        \item Legenda
        \item Indeks
        \item Skala
        \item Orientasi
        \item Inset
    \end{enumerate}

\item Diketahui jarak antara kota A dan kota B pada peta adalah 4~cm. Jika skala peta tersebut 1 : 1.000.000, maka jarak sebenarnya adalah \dots
    \begin{enumerate}[label=\alph*.]
        \item 4 km
        \item 40 km
        \item 400 km
        \item 4.000 km
        \item 40.000 km
    \end{enumerate}

\item Jika jarak sebenarnya dua kota adalah 150~km dan jarak pada peta 5~cm, berapakah skala peta tersebut?
    \begin{enumerate}[label=\alph*.]
        \item 1 : 3.000
        \item 1 : 30.000
        \item 1 : 300.000
        \item 1 : 3.000.000
        \item 1 : 30.000.000
    \end{enumerate}

\item Berdasarkan letak astronomisnya, garis bujur yang menjadi patokan Waktu Indonesia Barat (WIB) adalah \dots
    \begin{enumerate}[label=\alph*.]
        \item $105^{\circ}$ BT
        \item $120^{\circ}$ BT
        \item $135^{\circ}$ BT
        \item $95^{\circ}$ BT
        \item $141^{\circ}$ BT
    \end{enumerate}

\item Wilayah Waktu Indonesia Tengah (WITA) memiliki selisih waktu dari Greenwich Mean Time (UTC) sebesar \dots
    \begin{enumerate}[label=\alph*.]
        \item +6 jam
        \item +7 jam
        \item +8 jam
        \item +9 jam
        \item +10 jam
    \end{enumerate}

\item Provinsi di pulau Kalimantan yang masuk ke dalam zona Waktu Indonesia Barat (WIB) adalah \dots
    \begin{enumerate}[label=\alph*.]
        \item Kalimantan Timur dan Utara
        \item Kalimantan Selatan dan Timur
        \item Kalimantan Tengah dan Selatan
        \item Kalimantan Barat dan Tengah
        \item Kalimantan Barat dan Selatan
    \end{enumerate}

\item Jumlah provinsi terbaru di Indonesia saat ini adalah \dots
    \begin{enumerate}[label=\alph*.]
        \item 34 Provinsi
        \item 35 Provinsi
        \item 36 Provinsi
        \item 37 Provinsi
        \item 38 Provinsi
    \end{enumerate}

\item Pemekaran provinsi terbaru di Indonesia difokuskan di wilayah \dots
    \begin{enumerate}[label=\alph*.]
        \item Pulau Sumatera
        \item Pulau Jawa
        \item Pulau Kalimantan
        \item Pulau Sulawesi
        \item Pulau Papua
    \end{enumerate}

\item Syarat utama terjadinya interaksi sosial adalah adanya \dots
    \begin{enumerate}[label=\alph*.]
        \item Konflik dan akomodasi
        \item Kontak sosial dan komunikasi
        \item Toleransi dan pluralitas
        \item Nilai dan norma
        \item Asimilasi dan akulturasi
    \end{enumerate}

\item Manusia tidak dapat hidup sendiri dan selalu membutuhkan orang lain. Oleh karena itu, manusia disebut sebagai makhluk \dots
    \begin{enumerate}[label=\alph*.]
        \item Ekonomi
        \item Homo sapiens
        \item Zoon politicon
        \item Individu
        \item Biologis
    \end{enumerate}

\item Bermain game online bersama atau saling membalas email termasuk ke dalam jenis kontak sosial \dots
    \begin{enumerate}[label=\alph*.]
        \item Primer langsung
        \item Sekunder tidak langsung
        \item Primer tidak langsung
        \item Sekunder langsung
        \item Fisik
    \end{enumerate}

\item Kemajemukan masyarakat yang terdiri atas berbagai macam suku, agama, ras, dan golongan profesi disebut \dots
    \begin{enumerate}[label=\alph*.]
        \item Homogenitas
        \item Pluralitas
        \item Stratifikasi
        \item Diferensiasi
        \item Mobilitas
    \end{enumerate}

\item Sesuatu yang dianggap baik, berharga, pantas, dan dicita-citakan oleh masyarakat disebut \dots
    \begin{enumerate}[label=\alph*.]
        \item Norma sosial
        \item Nilai sosial
        \item Hukum adat
        \item Kebiasaan
        \item Interaksi sosial
    \end{enumerate}

\item Aturan yang bersumber dari hati nurani manusia mengenai apa yang benar dan salah disebut norma \dots
    \begin{enumerate}[label=\alph*.]
        \item Agama
        \item Kesusilaan
        \item Kesopanan
        \item Hukum
        \item Adat
    \end{enumerate}

\item Sanksi berupa rasa bersalah, penyesalan, atau malu merupakan sanksi dari pelanggaran norma \dots
    \begin{enumerate}[label=\alph*.]
        \item Hukum
        \item Agama
        \item Kesopanan
        \item Kesusilaan
        \item Kebiasaan
    \end{enumerate}

\item Menggunakan tangan kanan saat memberi sesuatu kepada orang lain merupakan wujud penerapan norma \dots
    \begin{enumerate}[label=\alph*.]
        \item Agama
        \item Kesusilaan
        \item Kesopanan
        \item Hukum
        \item Pidana
    \end{enumerate}

\item Aturan tertulis yang dibuat oleh lembaga negara yang berwenang dan bersifat memaksa disebut norma \dots
    \begin{enumerate}[label=\alph*.]
        \item Kebiasaan
        \item Kesusilaan
        \item Kesopanan
        \item Hukum
        \item Agama
    \end{enumerate}

\item Sanksi denda atau penjara diberikan kepada pelanggar norma \dots
    \begin{enumerate}[label=\alph*.]
        \item Hukum
        \item Kesopanan
        \item Agama
        \item Kesusilaan
        \item Sekolah
    \end{enumerate}

\item Tindakan cyberbullying di media sosial atau pencurian fasilitas sekolah termasuk ke dalam pelanggaran \dots
    \begin{enumerate}[label=\alph*.]
        \item Kedisiplinan
        \item Etika
        \item Kesopanan
        \item Hukum/Pidana
        \item Kesusilaan
    \end{enumerate}

\item Kondisi di mana hubungan sosial antaranggota masyarakat berlangsung stabil sesuai dengan nilai dan norma disebut \dots
    \begin{enumerate}[label=\alph*.]
        \item Mobilitas sosial
        \item Keteraturan sosial
        \item Dinamika sosial
        \item Konflik sosial
        \item Perubahan sosial
    \end{enumerate}

\item Proses penanaman nilai dan norma dari satu generasi ke generasi berikutnya disebut \dots
    \begin{enumerate}[label=\alph*.]
        \item Akomodasi
        \item Asimilasi
        \item Sosialisasi
        \item Mediasi
        \item Arbitrasi
    \end{enumerate}

\item Perpindahan posisi atau status sosial seseorang dari satu lapisan ke lapisan lain disebut \dots
    \begin{enumerate}[label=\alph*.]
        \item Interaksi sosial
        \item Keteraturan sosial
        \item Mobilitas sosial
        \item Konflik sosial
        \item Diferensiasi sosial
    \end{enumerate}

\item Seorang karyawan junior di sebuah perusahaan yang berkat prestasi kerjanya berhasil dipromosikan menjadi manajer. Ini adalah contoh mobilitas \dots
    \begin{enumerate}[label=\alph*.]
        \item Vertikal turun
        \item Horizontal
        \item Geografis
        \item Antargenerasi
        \item Vertikal naik
    \end{enumerate}

\item Seorang guru SMK pindah mengajar dari sekolah A ke sekolah B dengan jabatan yang sama. Hal ini merupakan mobilitas \dots
    \begin{enumerate}[label=\alph*.]
        \item Vertikal naik
        \item Vertikal turun
        \item Horizontal
        \item Statis
        \item Geografis
    \end{enumerate}

\item Perkembangan teknologi komunikasi dari pengiriman surat menjadi pesan instan (chatting) adalah bentuk dari \dots
    \begin{enumerate}[label=\alph*.]
        \item Konflik sosial
        \item Mobilitas sosial
        \item Keteraturan sosial
        \item Perubahan sosial dan budaya
        \item Sosialisasi
    \end{enumerate}

\item Suatu proses sosial ketika pihak-pihak berusaha memenuhi tujuannya dengan menentang lawan disertai ancaman atau kekerasan disebut \dots
    \begin{enumerate}[label=\alph*.]
        \item Persaingan
        \item Kontravensi
        \item Konflik
        \item Akomodasi
        \item Asimilasi
    \end{enumerate}

\item Proses penyesuaian diri untuk meredakan ketegangan tanpa menghancurkan pihak lawan disebut \dots
    \begin{enumerate}[label=\alph*.]
        \item Kompetisi
        \item Konflik
        \item Akomodasi
        \item Evolusi
        \item Revolusi
    \end{enumerate}

\item Bentuk akomodasi di mana masing-masing pihak yang berkonflik saling mengurangi tuntutannya demi mencapai kesepakatan disebut \dots
    \begin{enumerate}[label=\alph*.]
        \item Mediasi
        \item Kompromi
        \item Arbitrasi
        \item Konsiliasi
        \item Ajudikasi
    \end{enumerate}

\item Penyelesaian konflik yang melibatkan pihak ketiga sebagai penengah yang hanya memberikan nasihat dan tidak berhak memutuskan disebut \dots
    \begin{enumerate}[label=\alph*.]
        \item Arbitrasi
        \item Mediasi
        \item Kompromi
        \item Ajudikasi
        \item Toleransi
    \end{enumerate}

\item Penyelesaian masalah melalui jalur hukum di pengadilan disebut \dots
    \begin{enumerate}[label=\alph*.]
        \item Konsiliasi
        \item Arbitrasi
        \item Mediasi
        \item Ajudikasi
        \item Kompromi
    \end{enumerate}

\item Ilmu ekonomi pada dasarnya mempelajari bagaimana manusia memilih cara menggunakan sumber daya untuk \dots
    \begin{enumerate}[label=\alph*.]
        \item Menambah utang negara
        \item Mencari kekayaan sebanyak-banyaknya
        \item Memenuhi kebutuhan yang tidak terbatas
        \item Menghabiskan uang
        \item Menimbun barang berharga
    \end{enumerate}

\item Karakteristik dasar dari kebutuhan manusia adalah \dots
    \begin{enumerate}[label=\alph*.]
        \item Terbatas dan statis
        \item Tidak terbatas dan dinamis
        \item Dapat ditunda sepenuhnya
        \item Selalu sama setiap zaman
        \item Sangat mudah dipenuhi
    \end{enumerate}

\item Jurang pemisah antara kebutuhan yang tak terbatas dengan sumber daya yang terbatas menghasilkan masalah ekonomi utama yang disebut \dots
    \begin{enumerate}[label=\alph*.]
        \item Kelangkaan (Scarcity)
        \item Inflasi
        \item Deflasi
        \item Monopoli
        \item Kemiskinan
    \end{enumerate}

\item Kebutuhan mutlak yang harus dipenuhi pertama kali agar kelangsungan hidup manusia tidak terancam disebut kebutuhan \dots
    \begin{enumerate}[label=\alph*.]
        \item Sekunder
        \item Tersier
        \item Primer
        \item Imateril
        \item Mewah
    \end{enumerate}

\item Makanan, pakaian (sandang), dan tempat tinggal (papan) merupakan contoh dari kebutuhan \dots
    \begin{enumerate}[label=\alph*.]
        \item Tersier
        \item Primer
        \item Sekunder
        \item Rohani
        \item Imateril
    \end{enumerate}

\item Televisi, kasur empuk, dan kulkas biasanya digolongkan sebagai kebutuhan \dots
    \begin{enumerate}[label=\alph*.]
        \item Primer
        \item Sekunder
        \item Tersier
        \item Mewah
        \item Mutlak
    \end{enumerate}

\item Kebutuhan yang dipenuhi untuk menaikkan prestise atau status sosial seperti mobil sport dan berlian adalah kebutuhan \dots
    \begin{enumerate}[label=\alph*.]
        \item Primer
        \item Sekunder
        \item Tersier
        \item Pokok
        \item Jasmani
    \end{enumerate}

\item Berdasarkan sifatnya, kebutuhan manusia dibagi menjadi dua, yaitu \dots
    \begin{enumerate}[label=\alph*.]
        \item Primer dan Sekunder
        \item Individu dan Kelompok
        \item Sekarang dan Masa Depan
        \item Materil dan Imateril
        \item Pokok dan Tambahan
    \end{enumerate}

\item Buku pelajaran, nasi, dan pakaian merupakan contoh dari pemenuhan kebutuhan \dots
    \begin{enumerate}[label=\alph*.]
        \item Imateril
        \item Tersier
        \item Mewah
        \item Materil (Jasmani)
        \item Rohani
    \end{enumerate}

\item Kebutuhan yang menyangkut ketenangan jiwa, batin, dan emosional seperti beribadah dan hiburan disebut kebutuhan \dots
    \begin{enumerate}[label=\alph*.]
        \item Jasmani
        \item Materil
        \item Fisik
        \item Imateril (Rohani)
        \item Primer
    \end{enumerate}

\item Pelaku ekonomi yang memiliki peran ganda sebagai konsumen sekaligus pemilik faktor produksi adalah \dots
    \begin{enumerate}[label=\alph*.]
        \item Rumah Tangga Negara
        \item Masyarakat Luar Negeri
        \item Rumah Tangga Konsumsi (RTK)
        \item Rumah Tangga Produksi (RTP)
        \item Pemerintah
    \end{enumerate}

\item Imbalan yang diterima oleh Rumah Tangga Konsumsi (RTK) karena menyewakan atau menjual faktor produksi dapat berupa \dots
    \begin{enumerate}[label=\alph*.]
        \item Pajak dan retribusi
        \item Laba operasional perusahaan
        \item Upah, uang sewa, atau bunga
        \item Barang ekspor
        \item Subsidi pemerintah
    \end{enumerate}

\item Organisasi ekonomi yang bertujuan mengolah faktor produksi menjadi barang/jasa dan menjualnya untuk mendapat laba adalah \dots
    \begin{enumerate}[label=\alph*.]
        \item Rumah Tangga Konsumsi (RTK)
        \item Rumah Tangga Negara (RTN)
        \item Rumah Tangga Produksi (RTP)
        \item Lembaga Swadaya Masyarakat
        \item Koperasi simpan pinjam
    \end{enumerate}

\item Rumah Tangga Produksi (RTP) dituntut untuk selalu menerapkan prinsip efisiensi dan inovasi agar \dots
    \begin{enumerate}[label=\alph*.]
        \item Produknya laku di pasar dan mendapat laba
        \item Bisa memonopoli seluruh pasar
        \item Mendapat subsidi dari negara
        \item Tidak perlu membayar pajak
        \item Bisa mengekspor tanpa batas
    \end{enumerate}

\item Pelaku ekonomi yang memiliki kekuasaan hukum untuk membuat undang-undang, menetapkan tarif pajak, dan menjaga kestabilan harga adalah \dots
    \begin{enumerate}[label=\alph*.]
        \item Masyarakat Luar Negeri
        \item Rumah Tangga Konsumsi (RTK)
        \item Rumah Tangga Produksi (RTP)
        \item Rumah Tangga Negara (RTN)
        \item Investor asing
    \end{enumerate}

\item PT PLN dan Pertamina adalah contoh keterlibatan pemerintah (RTN) dalam perekonomian sebagai \dots
    \begin{enumerate}[label=\alph*.]
        \item Regulator
        \item Konsumen
        \item Distributor
        \item Produsen
        \item Pengawas
    \end{enumerate}

\item Ketika pemerintah membeli kendaraan dinas atau alat tulis kantor, maka pemerintah sedang berperan sebagai \dots
    \begin{enumerate}[label=\alph*.]
        \item Produsen
        \item Konsumen
        \item Regulator
        \item Distributor
        \item Kreditor
    \end{enumerate}

\item Perdagangan internasional berupa ekspor dan impor barang melibatkan peran pelaku ekonomi, yaitu \dots
    \begin{enumerate}[label=\alph*.]
        \item Rumah Tangga Konsumsi
        \item Rumah Tangga Produksi
        \item Masyarakat Luar Negeri
        \item Pemerintah Daerah
        \item Koperasi
    \end{enumerate}

\item Tujuan dari kegiatan ekspor barang ke luar negeri yang dilakukan oleh suatu negara adalah untuk \dots
    \begin{enumerate}[label=\alph*.]
        \item Memenuhi barang yang tidak bisa diproduksi
        \item Mendatangkan devisa
        \item Menurunkan nilai mata uang
        \item Mengurangi utang luar negeri
        \item Menambah angka kemiskinan
    \end{enumerate}

\item Mengimpor barang dari luar negeri biasanya dilakukan ketika \dots
    \begin{enumerate}[label=\alph*.]
        \item Negara sedang krisis moneter
        \item Masyarakat kelebihan produksi
        \item Kebutuhan tersebut tidak dapat diproduksi di dalam negeri
        \item Pemerintah ingin menghabiskan devisa
        \item Pajak impor sedang naik
    \end{enumerate}

\end{enumerate}
\end{document}